% Modelos de estrutura para o agente. Regras e fontes: ../AGENTS.md,
% seção "Referências e citações". Campos abaixo são marcadores, não obras reais.
% Copie o modelo pertinente, remova o % e substitua os campos por dados
% conferidos. Ordene os itens alfabeticamente. Não publique marcadores.
% Deixe uma linha vazia entre referências; main.tex controla o espaçamento.
% Não acrescente \\ no fim dos itens. Escape & como \& e _ como \_ no texto;
% dentro de \url{...}, mantenha o endereço verificado.
%
% LIVRO (edição e tradução somente quando pertinentes à edição consultada)
% SOBRENOME, Primeiro nome. \textbf{Título do livro}: subtítulo. Tradução de Nome do tradutor. 2. ed. Cidade: Editora, ANO.
%
% CAPÍTULO EM LIVRO ORGANIZADO (destaque no livro, não no capítulo)
% SOBRENOME, Primeiro nome. Título do capítulo. In: SOBRENOME, Primeiro nome (org.). \textbf{Título do livro}: subtítulo. Cidade: Editora, ANO. p. INICIAL--FINAL.
%
% ARTIGO ONLINE (destaque no periódico, não no artigo)
% SOBRENOME, Primeiro nome; SOBRENOME, Primeiro nome. Título do artigo. \textbf{Nome do periódico}, Cidade, v. VOLUME, n. NÚMERO, p. INICIAL--FINAL, ANO. DOI: \url{https://doi.org/DOI-VERIFICADO}. Disponível em: \url{URL-VERIFICADA}. Acesso em: DIA MÊS-ABREVIADO ANO.
% Se a publicação usa identificador eletrônico, substitua "p. INICIAL--FINAL"
% pelo identificador real (por exemplo, eNNNNN), sem inventar páginas ou número.
% DOI é incluído se houver. Volume, número e outros campos seguem a publicação.
%
% QUATRO OU MAIS AUTORES (convenção do projeto; a norma permite listar todos)
% SOBRENOME, Primeiro nome \textit{et al.}. Título do artigo. \textbf{Nome do periódico}, Cidade, v. VOLUME, n. NÚMERO, p. INICIAL--FINAL, ANO.
%
% TRABALHO EM EVENTO (destaque no título dos anais)
% SOBRENOME, Primeiro nome. Título do trabalho. In: NOME DO EVENTO, NÚMERO., ANO, Cidade do evento. \textbf{Anais [...]}. Cidade de publicação: Editora, ANO. p. INICIAL--FINAL.
% Adicione acesso online quando aplicável; confira dados do evento e dos anais.
%
% DISSERTAÇÃO OU TESE (confira grau, área, unidade e instituição na obra)
% SOBRENOME, Primeiro nome. \textbf{Título}: subtítulo. ANO. Dissertação (Mestrado em ÁREA) -- UNIDADE, INSTITUIÇÃO, Cidade, ANO.
% Para tese, adapte o tipo e o grau. Para consulta online, acrescente acesso.
%
% DOCUMENTO DE AUTORIA INSTITUCIONAL ONLINE
% NOME DA ENTIDADE. \textbf{Título do documento}: subtítulo. Cidade: Editora, ANO. Disponível em: \url{URL-VERIFICADA}. Acesso em: DIA MÊS-ABREVIADO ANO.
%
% PÁGINA WEB (confira autor responsável e data; o site não é autor automático)
% SOBRENOME, Primeiro nome. \textbf{Título da página}. Nome do site, DATA-DE-PUBLICAÇÃO. Disponível em: \url{URL-VERIFICADA}. Acesso em: DIA MÊS-ABREVIADO ANO.
%
% Dados ausentes: investigue e registre pendências no catálogo. Estes modelos
% não autorizam inventar dados nem usar marcadores como referências finais.

Nenhuma obra externa foi fornecida ou citada neste exemplo didático.
